\section{IMU Based Postural Stabilization}

Inertial corrections are computed asynchronously within the
strictly isolated Balance Controller. Raw quaternion orientation
data is converted to Euler roll ($\phi$) and pitch ($\theta$)
angles \cite{madgwick2011estimation} and processed through a
multistage transformation pipeline.

\subsection{Static Home Posture Correction}

During the static home phase, all four limbs maintain sustained
ground contact. The correction pipeline proceeds in five
sequential stages:

\textbf{Stage 1: EMA Prefiltering.}
Raw inertial measurements are smoothed via an exponential moving
average of the first order to suppress
sensor noise at high frequencies:
\begin{align}
    \hat{m}_k &= \alpha_{raw} \hat{m}_{k-1}
    + (1 - \alpha_{raw}) m_k \nonumber\\
    \alpha_{raw} &= 0.4
\end{align}

\textbf{Stage 2: Deadband Suppression.}
A configurable deadband of $\pm 0.02$ rad eliminates trivial
oscillations that would otherwise generate unnecessary servo
activity:
\begin{equation}
    e_{dz} =
    \begin{cases}
      e - d & e > d \\
      0 & |e| \leq d \\
      e + d & e < -d
    \end{cases}
    , \quad d = 0.02
\end{equation}

% Fig.~\ref{fig:deadband} illustrates the deadband transfer
% function, showing the suppression zone around the origin.
% 
% \begin{figure}[htbp]
%     \centering
%     \includegraphics[width=0.75\columnwidth]{figures/deadband_function.png}
%     \caption{Deadband suppression transfer function. Errors
%     within $\pm 0.02$ rad (red shaded region) are mapped to
%     zero, eliminating trivial oscillations. The dashed line
%     represents the unity response without deadband.}
%     \label{fig:deadband}
% \end{figure}

\textbf{Stage 3: Gain Scheduled PID with Antiwindup.}
A full PID controller \cite{astrom2006advanced} with
three novel enhancements processes the deadband filtered error:
\begin{align}
    u_{PID} &= K_P \cdot g_s \cdot e \nonumber\\
    &+ K_I \cdot g_s \cdot \mathcal{I}
    + K_D \cdot \dot{e}_{lpf}
\end{align}

\noindent where the gain schedule factor $g_s \in [0.5, 1.0]$
scales proportionally to the error magnitude, preventing
corrections that are overly aggressive near zero:
\begin{equation}
    g_s = \text{clamp}\!\left(
    \frac{|e|}{e_{thresh}},\, 0.5,\, 1.0
    \right)
\end{equation}

\noindent where $e_{thresh} = 0.08$ rad.
% Fig.~\ref{fig:gain_schedule} visualizes this piecewise linear
% relationship.
% 
% \begin{figure}[htbp]
%     \centering
%     \includegraphics[width=0.75\columnwidth]{figures/gain_schedule.png}
%     \caption{Proportional gain scheduling factor $g_s$ as a
%     function of error magnitude. Below $|e|=0.04$ rad, $g_s$
%     is clamped at 0.5; above $e_{thresh}=0.08$ rad, it
%     saturates at 1.0.}
%     \label{fig:gain_schedule}
% \end{figure}

The integral accumulator features exponential leakage decay
($\lambda = 0.9995$) to prevent windup at steady state, bounded
by symmetric saturation limits ($|\mathcal{I}| \leq 1.5$).
Additionally, back calculation antiwindup reduces the
accumulator when the PID output exceeds saturation bounds:
\begin{equation}
    \mathcal{I}_{k} = \mathcal{I}_{k} -
    \frac{u_{raw} - u_{sat}}{2 K_I}
    \quad \text{if} \quad |u_{raw}| > u_{max}
\end{equation}

The derivative term is presmoothed with a dedicated LPF
($\alpha_d = 0.7$) and includes a derivative ramp threshold
($e_{ramp} = 0.005$) that zeroes the derivative when the error
is negligibly small, preventing derivative kick.

\textbf{Stage 4: Output Saturation.}
The composite PID output is clamped to
$|u| \leq 1.5$ rad to protect actuator limits.

\textbf{Stage 5: Output Lowpass Filtering.}
A final lowpass filter of the first order smooths the corrective
output to eliminate residual step discontinuities:
\begin{align}
    H_{k} &= \alpha_c H_{k-1}
    + (1 - \alpha_c)\, u_{sat} \nonumber\\
    \alpha_c &= 0.15
\end{align}

\subsection{Phase Aware Dynamic Gait Correction}

While the robot is actively ambulating, the correction pipeline
switches to an adaptive PID controller specific to the gait cycle
that incorporates three additional mechanisms:

\textbf{Moving Average Baseline.}
IMU measurements are accumulated into a buffer spanning one full
gait cycle. The running mean provides a smoothed estimate of the
true tilt over the cycle, filtering out periodic oscillations
inherent to legged locomotion.

\textbf{Adaptive Proportional Gain.}
The proportional gain scales dynamically with the instantaneous
tilt magnitude to provide stronger correction authority during
large disturbances:
\begin{equation}
    K_{P,eff} = K_P \left(1 + \min\!\left(1.0, \;
    \frac{\max(|\phi|, |\theta|)}{K_{prop}}
    \right)\right)
\end{equation}

\textbf{Zero Crossing Integral Reset.}
The integral accumulator is immediately reset to zero whenever
the error signal crosses zero, preventing oscillatory overshoot
during transient recovery.

\textbf{Stance Phase Rotational Compensation.}
The corrective output is applied as a rotation matrix in the
body frame $R = R_y(\theta) R_x(\phi)$:
\begin{equation}
    R = \begin{bmatrix}
    c_\theta & s_\phi s_\theta & c_\phi s_\theta \\
    0 & c_\phi & -s_\phi \\
    -s_\theta & s_\phi c_\theta & c_\phi c_\theta
    \end{bmatrix}
\end{equation}

This matrix is applied exclusively to limbs satisfying the
stance condition $|Z_{foot} - Z_{stance}| \leq Z_{thresh}$,
where $Z_{stance} = -170$ mm and $Z_{thresh} = 5$ mm. Only
the vertical component of the transformed foot position is
modified; horizontal coordinates are preserved to maintain
stride geometry:
\begin{align}
    \mathbf{p}_{adj} &= R \cdot \mathbf{p}_{raw} \\
    x_{adj} &= x_{raw}, \quad y_{adj} = y_{raw} \nonumber
\end{align}
